% !TEX root = ../main.tex
\section{Proof of Theorem~\ref{thm:upper-beta1}}
\label{app:upper-beta1}

\subsection{Curvature and smoothness of the potential function}
\label{subsec:curvature-smoothness}
Let us define the constants
\begin{equation}
    L_{\calD,\bx}
    :=
    L_f + L_{\calA}L_h + \rho G_{\calA}^{2},
    \label{eq:x-weak-convexity-modulus}
\end{equation}
and 
\[
    \mu_\gamma
    :=
    \frac{1}{\gamma}-L_{\calD,\bx},
    \qquad
    L_x
    :=
    L_{\calD,\bx}+\frac{1}{\gamma},
    \qquad
    L_\lambda(s)
    :=
    \frac{1}{\rho}+\frac{1}{s},
\] where $L_f$, $L_{\calA}$, $G_{\calA}$, and $L_h$ are the constants of
Assumptions~\ref{ass:smoothness} and~\ref{ass:h}. 

The following proposition summarizes the curvature and smoothness properties
of the smoothed perturbed potential function $\mathcal P_s$.
The proof is deferred to Appendix~\ref{app:curvature-smoothness}.

\begin{proposition}
\label{prop:curvature-smoothness}
Suppose that Assumptions~\ref{ass:smoothness} and~\ref{ass:h} hold, and let the smoothing parameter  
$ 0 < \gamma < \frac{1}{2L_{\calD,\bx}}$. Then, for all $s>0$, $\mathcal P_s$ is
$\mu_\gamma $-strongly convex in $\bx$ and
$1/s$-strongly concave in $\blmbd$. Moreover, for all
$\bx,\bx'\in\RR^n$,
$\blmbd,\blmbd'\in\RR^m$,
$\bz\in\RR^n$, and
$\bmu\in\RR^m$, it holds that 
\begin{align}
    &\bigl\|
    \nabla_{\bx}\mathcal{P}_s(\bx,\blmbd,\bz,\bmu)
    -
    \nabla_{\bx}\mathcal{P}_s(\bx',\blmbd',\bz,\bmu)
    \bigr\|
    \le
    L_x\|\bx-\bx'\|
    +
    G_{\calA}\|\blmbd-\blmbd'\|,
    \label{eq:payoff-x-smoothness}
    \\
    &\bigl\|
    \nabla_{\blmbd}\mathcal{P}_s(\bx,\blmbd,\bz,\bmu)
    -
    \nabla_{\blmbd}\mathcal{P}_s(\bx',\blmbd',\bz,\bmu)
    \bigr\|
    \le
    G_{\calA}\|\bx-\bx'\|
    +
    L_\lambda(s)\|\blmbd-\blmbd'\|
    \label{eq:payoff-lambda-smoothness}.
\end{align}
In particular, $\mathcal P_s$ is $L_x$-smooth in $\bx$ and $L_\lambda(s)$-smooth in $\blmbd$.
\end{proposition}

\subsection{Proof of the iteration complexity bound}

Throughout this appendix, let
\[
L:=\sqrt2\,\max\Bigl\{L_{\calD,\bx},\ G_{\calA},\ \frac1\rho\Bigr\},
\qquad
\bar\gamma:=\frac1{4L},
\qquad
\bar\tau:=\frac1{36L}\Bigl(\frac\eta\gamma\Bigr)^2,
\]
and let \(\mathcal Y:=\{\blmbd\in\RR^m:\|\blmbd\|\le L_h+1\}\). We write
\([\mathrm{ZXSL}]\) for~\cite{zhang2020smoothed}; equation and lemma numbers
refer to its arXiv version~2010.15768v3.

\paragraph{Step 1: reduction to Smoothed GDA}
Since \(\beta=1\) and \(\bmu^0=\blmbd^0=\zero\), induction gives
\(\bmu^k=\blmbd^k\) for all \(k\). Hence the gradient of the dual regularizer
vanishes at every dual step, and the updates~\eqref{eq:splalm-updates} read
\[
\bx^{k+1}=\bx^k-\eta\Bigl[\nabla_{\bx}\calD_\rho(\bx^k,\blmbd^k)
  +\frac1\gamma(\bx^k-\bz^k)\Bigr],
\qquad
\blmbd^{k+1}=\blmbd^k+\tau\nabla_{\blmbd}\calD_\rho(\bx^{k+1},\blmbd^k),
\]
together with \(\bz^{k+1}=\bz^k+\alpha(\bx^{k+1}-\bz^k)\). None of these
updates involves \(s\). By~\eqref{eq:prox-multiplier},
\(\nabla_{\blmbd}\calD_\rho(\bx,\blmbd)=\calA(\bx)-\by(\bx,\blmbd)
=(\bq(\bx,\blmbd)-\blmbd)/\rho\), so
\[
\blmbd^{k+1}=\Bigl(1-\frac\tau\rho\Bigr)\blmbd^k+\frac\tau\rho\,\bq^k,
\qquad \|\bq^k\|\le L_h,
\]
where the bound holds because \(\bq^k\in\partial h(\by^{k+1})\) and \(h\) is
\(L_h\)-Lipschitz. Since \(L\ge1/\rho\), we have
\(\tau\le\bar\tau\le1/(144L)\le\rho\), and induction gives
\(\|\blmbd^k\|\le L_h\) for all \(k\). Hence every \(\blmbd^k\) lies in the
interior of \(\mathcal Y\), and the projection onto \(\mathcal Y\) does not
change the iterates. Consequently, \(\{(\bx^k,\blmbd^k,\bz^k)\}\) is the
sequence generated by Smoothed GDA~\cite[Algorithm~2]{zhang2020smoothed}
applied to
\(\min_{\bx\in\RR^n}\max_{\blmbd\in\mathcal Y}\calD_\rho(\bx,\blmbd)\) from
zero, with smoothing parameter \(p=1/\gamma\), primal stepsize \(\eta\), dual
stepsize \(\tau\), and averaging parameter \(\alpha\).

\paragraph{Step 2: assumptions of~\cite[Theorem~3.4]{zhang2020smoothed}}
The proof of Proposition~\ref{prop:curvature-smoothness} gives
\[
\begin{aligned}
\|\nabla_{\bx}\calD_\rho(\bx,\blmbd)-\nabla_{\bx}\calD_\rho(\bx',\blmbd')\|
&\le L_{\calD,\bx}\|\bx-\bx'\|+G_{\calA}\|\blmbd-\blmbd'\|,\\
\|\nabla_{\blmbd}\calD_\rho(\bx,\blmbd)-\nabla_{\blmbd}\calD_\rho(\bx',\blmbd')\|
&\le G_{\calA}\|\bx-\bx'\|+\frac1\rho\|\blmbd-\blmbd'\|.
\end{aligned}
\]
By the Cauchy--Schwarz inequality, both gradients are \(L\)-Lipschitz on
\(\RR^n\times\RR^m\). The set \(\mathcal Y\) is convex and compact,
\(\RR^n\) is closed and convex, and \(\calD_\rho(\bx,\cdot)\) is concave. By
Lemma~\ref{lem:effective-penalty},
\[
\calD_\rho(\bx,\blmbd)
=\max_{\bq}\Bigl\{\calL(\bx,\bq)-\frac1{2\rho}\|\bq-\blmbd\|^2\Bigr\}
\le\max_{\bq}\calL(\bx,\bq)=F(\bx),
\]
with equality for \(\blmbd\in\partial h(\calA(\bx))\). Every such
\(\blmbd\) satisfies \(\|\blmbd\|\le L_h\) and hence lies in \(\mathcal Y\).
Therefore \(\max_{\blmbd\in\mathcal Y}\calD_\rho(\bx,\blmbd)=F(\bx)\), which
is bounded below by \(\inf F>-\infty\). Thus Assumption~3.3 of
\([\mathrm{ZXSL}]\) holds with gradient Lipschitz constant \(L\).

\paragraph{Step 3: parameter conditions}
We verify condition~(3.3) of \([\mathrm{ZXSL}]\). First,
\(\gamma\le\bar\gamma\) gives \(p=1/\gamma\ge4L>3L\). Second, since
\(\gamma L\le1/4\),
\(\gamma/(1+\gamma L)\ge4\gamma/5>\gamma/2\ge\eta\), hence
\(\eta<1/(p+L)\). Third, put \(a:=\eta(p-L)\). Then
\(3\eta/(4\gamma)\le a<\eta p\le1/2\), so
\[
\frac{\eta^2(p-L)^2}{4L(1+\eta(p-L))^2}
>\frac{a^2}{16L}
\ge\frac9{256L}\Bigl(\frac\eta\gamma\Bigr)^2
>\bar\tau\ge\tau,
\]
and \(\bar\tau\le1/(144L)<1/(11L)\). Finally, the averaging parameter
\(\alpha=c\epsilon^2\le c\) satisfies
\(\alpha\le\min\{1/36,(p-L)^2/(384p(p+L)^2)\}\) by the choice of \(c\) in
Step~5.

\paragraph{Step 4: initial gap}
Let \(K(\bx,\bz;\blmbd):=\calD_\rho(\bx,\blmbd)+\frac1{2\gamma}\|\bx-\bz\|^2\),
\(d(\blmbd,\bz):=\min_{\bx}K(\bx,\bz;\blmbd)\), and
\(P(\bz):=\min_{\bx}\max_{\blmbd\in\mathcal Y}K(\bx,\bz;\blmbd)\). The
potential function of~\cite[(4.1)]{zhang2020smoothed} at the zero
initialization is \(\varphi^0=K(\zero,\zero;\zero)-2d(\zero,\zero)+2P(\zero)\).
By Proposition~\ref{prop:curvature-smoothness}, \(K(\cdot,\zero;\zero)\) is
\(\mu_\gamma\)-strongly convex, and
\(\|\nabla_{\bx}\calD_\rho(\zero,\zero)\|
=\|\nabla f(\zero)+J_{\calA}(\zero)^\top\bq(\zero,\zero)\|\le G_f+G_{\calA}L_h\).
Hence
\[
d(\zero,\zero)\ge\calD_\rho(\zero,\zero)-\frac{(G_f+G_{\calA}L_h)^2}{2\mu_\gamma}.
\]
By Step~2, \(P(\zero)=\min_{\bx}\{F(\bx)+\frac1{2\gamma}\|\bx\|^2\}\le F(\zero)\).
Moreover,
\(\calD_\rho(\zero,\zero)=f(\zero)+\min_{\by}\{h(\by)+\frac\rho2\|\calA(\zero)-\by\|^2\}\),
and the Lipschitz continuity of \(h\) gives
\(0\le F(\zero)-\calD_\rho(\zero,\zero)\le L_h^2/(2\rho)\). Combining these
bounds,
\[
\varphi^0-\inf F\le\hat\Delta
:=\Delta+\frac{L_h^2}{2\rho}+\frac{(G_f+G_{\calA}L_h)^2}{\mu_\gamma},
\]
where \(\mu_\gamma\ge1/\gamma-L\ge3/(4\gamma)>0\).

\paragraph{Step 5: rate}
For \(t\ge0\), set
\[
S_{t+1}:=\max\bigl\{\|\nabla_{\bx}\calD_\rho(\bx^{t+1},\blmbd^{t+1})\|,\
\|\nabla_{\blmbd}\calD_\rho(\bx^{t+1},\blmbd^{t+1})\|\bigr\}.
\]
Since \(\blmbd^{t+1}\) lies in the interior of \(\mathcal Y\), the normal
cone of \(\mathcal Y\) at \(\blmbd^{t+1}\) is \(\{\zero\}\), and
\((\bx^{t+1},\blmbd^{t+1})\) is an \(S_{t+1}\)-solution in the sense
of~\cite[Definition~3.1]{zhang2020smoothed}. Fix \(T\ge1\). The proof
of~\cite[Theorem~3.4]{zhang2020smoothed} distinguishes the cases (B.52) and
(B.53) and yields some \(t\in\{0,\dots,T-1\}\) with either
\[
\text{(i)}\quad S_{t+1}\le\lambda_1\max\{\alpha,\sqrt\alpha\}
\qquad\text{or}\qquad
\text{(ii)}\quad S_{t+1}\le\lambda_2\sqrt{\frac{\varphi^0-\inf F}{T\alpha}},
\]
by (B.62)--(B.63) with Lemma~B.11 in case (i) and by Lemma~B.12 in case (ii).
Here \(\lambda_1\) and \(\lambda_2\) depend only on \(L\), \(p\), \(\eta\),
\(\tau\), and the diameter \(2(L_h+1)\) of \(\mathcal Y\). Choose
\[
c:=\min\Bigl\{\frac1{36},\ \frac{(p-L)^2}{384p(p+L)^2},\
\Bigl(\frac{10}{11\lambda_1}\Bigr)^2\Bigr\},
\qquad
T:=\Bigl\lceil\Bigl(\frac{11}{10}\Bigr)^2\frac{\lambda_2^2\hat\Delta}{c\,\epsilon^4}\Bigr\rceil.
\]
Since \(\alpha=c\epsilon^2\le1\), case (i) gives
\(S_{t+1}\le\lambda_1\sqrt c\,\epsilon\le\frac{10}{11}\epsilon\), and case
(ii) gives \(S_{t+1}\le\frac{10}{11}\epsilon\) by Step~4 and the choice of
\(T\).

\paragraph{Step 6: from \(S_{t+1}\) to the KKT residual}
By~\eqref{eq:kkt-residual} and~\eqref{eq:prox-multiplier},
\(\mathcal R_{\mathrm{KKT}}(\bx^{t+1},\by^{t+1},\bq^t)
=\max\{\|\nabla_{\bx}\calD_\rho(\bx^{t+1},\blmbd^t)\|,\|\mathbf g\|\}\), where
\(\mathbf g:=\nabla_{\blmbd}\calD_\rho(\bx^{t+1},\blmbd^t)\) and
\(\blmbd^{t+1}-\blmbd^t=\tau\mathbf g\). By Step~2 and \(\tau/\rho\le1/144\),
\(\|\nabla_{\blmbd}\calD_\rho(\bx^{t+1},\blmbd^{t+1})-\mathbf g\|\le\|\mathbf g\|/144\),
hence \(\|\mathbf g\|\le\frac{144}{143}S_{t+1}\). Similarly,
\(G_{\calA}\tau\le1/144\) gives
\(\|\nabla_{\bx}\calD_\rho(\bx^{t+1},\blmbd^t)\|
\le S_{t+1}+G_{\calA}\tau\|\mathbf g\|\le\frac{144}{143}S_{t+1}\). Therefore
\(\mathcal R_{\mathrm{KKT}}(\bx^{t+1},\by^{t+1},\bq^t)\le\frac{11}{10}S_{t+1}\le\epsilon\),
and
\[
T^{\mathrm{KKT}}_\epsilon(\theta,\alpha,1;I)\le t+1\le T\le C\epsilon^{-4},
\qquad
C:=\Bigl(\frac{11}{10}\Bigr)^2\frac{\lambda_2^2\hat\Delta}{c}+1,
\]
where the last inequality uses \(\epsilon\le1\). The constant \(c\) depends
only on \(L\), \(\gamma\), \(\eta\), \(\tau\), and \(L_h\), hence only on
\(\theta\) and \(\mathbf C\); the constant \(C\) depends in addition on
\(\Delta\) through \(\hat\Delta\). Neither depends on the instance or on the
dimensions.
\hfill\(\square\)
